% ---------------------------------------------------------------
% کامپایل با LuaLaTeX:   lualatex main_fa.tex   (دو بار برای فهرست مطالب)
% فونت‌های لازم: B Nazanin (معمولی و Bold) و Courier New
% ---------------------------------------------------------------
\documentclass[12pt,a4paper]{article}

\usepackage[no-math]{fontspec}
\usepackage[bidi=basic,english]{babel}
\babelprovide[import,main,maparabic,mapdigits]{persian}

\babelfont[persian]{rm}[Scale=1.15]{B Nazanin}
\babelfont[persian]{sf}[Scale=1.15]{B Nazanin}
\babelfont[persian]{tt}{Courier New}
\babelfont[english]{rm}[Scale=0.833]{Times New Roman}
\babelfont[english]{sf}[Scale=0.833]{Times New Roman}
\babelfont[english]{tt}{Courier New}

% B Nazanin lacks U+0654; use its supported precomposed ezafe character U+06C0.
\usepackage[margin=0.85in,headheight=14pt]{geometry}
\usepackage{amsmath,amssymb}
\usepackage{graphicx}
\usepackage{xcolor}
\usepackage{booktabs}
\usepackage{tabularx}
\usepackage{array}
\usepackage{enumitem}
\usepackage{listings}
\usepackage{tikz}
\usetikzlibrary{arrows.meta,positioning,fit,calc,shapes.geometric}
\usepackage{pgfplots}
\pgfplotsset{compat=1.18}
\usepackage{fancyhdr}
\usepackage{hyperref}
\usepackage[nameinlink,noabbrev]{cleveref}

% نام‌های فارسی برای ارجاع‌ها
\crefname{figure}{شکل}{شکل‌ها}
\Crefname{figure}{شکل}{شکل‌ها}
\crefname{table}{جدول}{جدول‌ها}
\Crefname{table}{جدول}{جدول‌ها}
\crefname{section}{بخش}{بخش‌ها}
\Crefname{section}{بخش}{بخش‌ها}

% متن لاتین داخل متن فارسی
\providecommand{\lr}[1]{\foreignlanguage{english}{\begingroup\hyphenchar\font=-1 #1\endgroup}}
\newcommand{\fa}[1]{\foreignlanguage{persian}{#1}}
\newcommand{\code}[1]{\lr{\texttt{#1}}}

\definecolor{primaryblue}{RGB}{30,58,138}
\definecolor{midblue}{RGB}{59,130,246}
\definecolor{lightblue}{RGB}{239,246,255}
\definecolor{darkgray}{RGB}{51,65,85}
\definecolor{lightgray}{RGB}{248,250,252}
\definecolor{teal}{RGB}{13,148,136}
\definecolor{orange}{RGB}{234,88,12}
\definecolor{redtone}{RGB}{185,28,28}

\hypersetup{
    colorlinks=true,
    linkcolor=primaryblue,
    citecolor=primaryblue,
    urlcolor=midblue,
    pdftitle={SonicDrop: Acoustic Data Communication},
    pdfauthor={Reza Ahmadizadeh}
}

\pagestyle{fancy}
\fancyhf{}
\lhead{\small \lr{SonicDrop}: ارتباط داده با صوت}
\rhead{\small گزارش فنی پروژه}
\cfoot{\thepage}
\renewcommand{\arraystretch}{1.25}
\setlength{\parskip}{0.45em}
\setlength{\parindent}{0pt}
\setlist[itemize]{label=\ensuremath{\bullet},leftmargin=1.5em,itemsep=0.2em,topsep=0.25em}
\setlist[enumerate]{leftmargin=1.8em,itemsep=0.25em,topsep=0.25em}

\tikzset{
    block/.style={draw=primaryblue,fill=lightblue,rounded corners=2pt,
        align=center,minimum height=0.9cm,text width=2.35cm,font=\small},
    titlebox/.style={draw=primaryblue,fill=lightblue,rounded corners=3pt,
        align=center,minimum width=3.9cm,minimum height=1.35cm,
        text width=3.6cm,font=\small},
    architecturebox/.style={draw=primaryblue,fill=lightblue,rounded corners=3pt,
        align=center,minimum width=2.8cm,minimum height=1.1cm,
        text width=2.6cm,font=\footnotesize},
    smallblock/.style={draw=darkgray,fill=lightgray,rounded corners=2pt,
        align=center,minimum height=0.75cm,text width=1.65cm,font=\footnotesize},
    arrow/.style={-{Stealth[length=2mm]},thick,draw=darkgray},
    dashedarrow/.style={-{Stealth[length=2mm]},thick,dashed,draw=orange},
    node distance=0.55cm and 0.45cm,
    every node/.append style={execute at begin node=\selectlanguage{persian}}
}

\begin{document}

\begin{titlepage}
\centering
\vspace*{1.6cm}
{\Large\bfseries گزارش فنی پروژه\par}
\vspace{0.8cm}
{\Huge\bfseries\color{primaryblue} \lr{SonicDrop}\par}
\vspace{0.35cm}
{\LARGE فرستنده‌گیرندۀ نرم‌افزارمحورِ دادۀ صوتی\par}
\vspace{0.6cm}
{\large انتقال مطمئن متن و تصویر دودویی روی کانال صوتی\par}
\vspace{1.8cm}

% جریان نمودار از راست به چپ (متناسب با جهت خواندن فارسی)
\begin{otherlanguage}{english}
\begin{tikzpicture}[>=Stealth]
    \node[titlebox] (input) at (10.45,0) {داده ورودی\\متن / فکس $16\times16$};
    \node[titlebox] (source) at (6.15,0) {کدگذاری منبع\\\lr{UTF-8} / \lr{RLE} اختیاری};
    \node[titlebox] (packet) at (1.85,0) {حفاظت بسته\\\lr{CRC / FEC / HDLC}};
    \node[titlebox] (modem) at (1.85,-2.5) {شکل‌موج \lr{M-FSK}\\صدا / \lr{WAV}};
    \node[titlebox] (channel) at (6.15,-2.5) {کانال صوتی\\بلندگو، هوا، میکروفون};
    \node[titlebox] (receiver) at (10.45,-2.5) {کدگشایی و بازسازی\\همگام‌سازی، متن، فکس};
    \draw[arrow] (input)--(source);
    \draw[arrow] (source)--(packet);
    \draw[arrow] (packet)--(modem);
    \draw[arrow] (modem)--(channel);
    \draw[arrow] (channel)--(receiver);
\end{tikzpicture}
\end{otherlanguage}

\vfill
\begin{tabular}{lr}
    \lr{Reza Ahmadizadeh} & \textbf{نویسنده:} \\
    مودم صوتی (\lr{SonicDrop}) & \textbf{پروژه:} \\
    \lr{Python}، \lr{Django}، \lr{NumPy} و \lr{API}‌های صوتی مرورگر & \textbf{پیاده‌سازی:} \\
    \today & \textbf{تاریخ گزارش:}
\end{tabular}
\vspace*{1.4cm}
\end{titlepage}

\begin{abstract}
\lr{SonicDrop} یک مودم نرم‌افزارمحور و برنامۀ وب برای ارسال پیام‌های متنی کوتاه و تصاویر دودویی کوچک از طریق صوت است. فرستنده داده‌های برنامه را به بسته‌ها تبدیل می‌کند، یک سرآیند فشرده و \lr{CRC-8} به آن‌ها می‌افزاید، به‌صورت اختیاری کدگذاری طول‌دنباله (\lr{RLE}) و تصحیح خطای پیشرو (\lr{FEC}) همینگ \lr{(7,4)} را اعمال می‌کند، رشتۀ بیت‌ها را با پرچم‌ها و درج بیت به سبک \lr{HDLC} قاب‌بندی می‌کند و یک شکل‌موج تک‌کانالۀ \lr{M-FSK} تولید می‌کند. گیرنده یک فایل \lr{WAV} یا صدای میکروفون را می‌پذیرد، تن‌های کاندیدا را آشکار می‌کند، چند انحراف زمانیِ نماد را جست‌وجو می‌کند، قاب‌ها را بررسی می‌کند، خطاهای بیتیِ قابل‌تصحیح را اصلاح می‌کند و محتوای بسته‌ها را بازسازی می‌کند. رابط کاربری همچنین نویز شبیه‌سازی‌شدۀ کانال، معیارهای سیگنال، بازرسی قاب و یک نمایش \lr{ACK/NACK} صوتی هماهنگ‌شده در مرورگر را ارائه می‌دهد. این گزارش خط‌لولۀ پیاده‌سازی‌شده و مفروضات آن را شرح می‌دهد، محاسبات اسمی را از نتایج اندازه‌گیری‌شده جدا می‌کند و محدودیت‌های مهم نمونۀ اولیه را مشخص می‌کند. تبدیل سفارشی گذرکلید، رمزنگاریِ مناسب برای استفاده در محیط تولید نیست. همچنین، نمایش \lr{ARQ} در مرورگر یک پروتکل انتقال کامل به‌شمار نمی‌آید.
\end{abstract}

\tableofcontents
\clearpage

\section{مقدمه}
\label{sec:introduction}

بیشتر دستگاه‌های مصرفی داده‌های خود را از طریق شبکه‌های رادیویی یا اتصال‌های فیزیکی مبادله می‌کنند. \lr{SonicDrop} کانال متفاوتی را بررسی می‌کند: بلندگو دنباله‌ای از تن‌های صوتی منتشر می‌کند و میکروفون آن‌ها را ضبط می‌کند. این طرح به‌عنوان یک نمونۀ اولیۀ آموزشی برای ارتباط موردی و کوتاه‌برد، آزمایش‌های پردازش سیگنال و نمایش بصریِ فنون لایۀ پیوند داده در نظر گرفته شده است. این روش نیازی ندارد که دو سر ارتباط \lr{Wi-Fi} یا بلوتوث مشترک داشته باشند، هرچند برنامۀ وب فعلی برای آماده‌سازی و بررسی ارسال‌ها از یک مرورگر و یک سرور \lr{Django} استفاده می‌کند.

اتاق آکوستیک محیط ارتباطی دشواری است. بازتاب‌ها می‌توانند نسخه‌های تأخیریافتۀ یک تن را بسازند که با نمادهای بعدی هم‌پوشانی پیدا می‌کنند؛ پاسخ بلندگو و میکروفون در سراسر فرکانس‌ها تخت نیست؛ فاصله و جهت‌گیری سطح سیگنال دریافتی را تغییر می‌دهند؛ و صداهای نامرتبط می‌توانند سیگنال را بپوشانند. \lr{SonicDrop} با این مسائل از طریق انتخاب‌های متعدد تن، لبه‌های نماد با شیب نرم، ناحیۀ انتگرال‌گیریِ کوتاه‌شده در گیرنده، نرمال‌سازی انرژی هر حامل، همگام‌سازی قاب، آشکارسازی خطا، \lr{FEC} اختیاری و مرزهای تکراری بسته مقابله می‌کند.

\subsection{اهداف پروژه}
پروژۀ پیاده‌سازی‌شده موارد زیر را نشان می‌دهد:
\begin{itemize}
    \item کدگذاری متن \lr{UTF-8} و یک تصویر دودویی $16\times16$ به‌صورت بیت.
    \item ارسال بیت‌ها با \lr{2-FSK} و \lr{4-FSK} در باند شنیداری یا نزدیک‌به‌فراصوت.
    \item حفظ مرز قاب‌ها و شناسایی قاب‌های آسیب‌دیده با استفاده از پرچم‌ها، درج بیت و \lr{CRC-8}.
    \item تصحیح یک خطا در هر کدواژۀ همینگ \lr{(7,4)} در صورت فعال بودن این گزینه.
    \item فشرده‌سازی اختیاری رشته‌های بیتی با یک قالب \lr{RLE} ساده و بدون اتلاف.
    \item تولید صدای \lr{WAV}، ضبط صدای میکروفون در مرورگر و رمزگشایی فایل‌های \lr{WAV} بارگذاری‌شده.
    \item نمایش شکل‌موج، قاب، \lr{BER}، \lr{SNR} و اطلاعات ظرفیت نظری کانال.
\end{itemize}

\subsection{دامنه و اصطلاحات}
واژه‌های \textbf{فرستنده} و \textbf{گیرنده} به نقش‌های سرتاسری (انتها‌به‌انتها) اشاره دارند. برنامۀ \lr{Django} یک ارسال را می‌سازد و یک فایل صوتی برمی‌گرداند؛ مرورگر می‌تواند آن را پخش کند یا ورودی را ضبط کند. کاربر همچنین می‌تواند فایل \lr{WAV} را به دستگاه دیگری منتقل کند و برای رمزگشایی بارگذاری کند. برخی رفتارهای رابط کاربری، به‌ویژه نمایش زندۀ بسته‌ها و هماهنگی \lr{ACK/NACK}، بین زمینه‌های مرورگر هماهنگ می‌شود؛ بنابراین باید آن را نمایشی دانست که پیرامون خط‌لولۀ صوتی ساخته شده است، نه شاهدی بر یک پروتکل شبکه‌ای کاملاً خودمختار.

\section{نمای کلی سیستم}
\label{sec:architecture}

\Cref{fig:architecture} مسیر منطقی را در پروژه نشان می‌دهد. در فرستنده، کدگذاری از راست به چپ پیش می‌رود و در گیرنده، رمزگشایی تبدیل‌ها را وارونه می‌کند. کانال میانی همان مسیر فیزیکی هواست و ممکن است نویز، تضعیف، انحراف زمانی و اثرات چندمسیری اضافه کند.

\begin{figure}[ht]
\centering
\begin{otherlanguage}{english}
\begin{tikzpicture}[>=Stealth]
    \node[architecturebox] (app) at (13.525,0) {داده برنامه\\متن / فکس};
    \node[architecturebox] (source) at (10.475,0) {کدگذاری منبع\\\lr{UTF-8 / RLE}};
    \node[architecturebox] (packet) at (7.425,0) {بسته‌سازی\\سرآیند ۱۲ بیتی};
    \node[architecturebox] (link) at (4.375,0) {\lr{HDLC + FEC}\\\lr{CRC} / همینگ};
    \node[architecturebox] (phy) at (1.325,0) {ارسال \lr{M-FSK}\\تن / \lr{WAV}};
    \node[architecturebox] (out) at (13.525,-2.15) {خروجی برنامه\\متن / فکس};
    \node[architecturebox] (sourceout) at (10.475,-2.15) {کدگشایی منبع\\\lr{RLE / UTF-8}};
    \node[architecturebox] (linkout) at (7.425,-2.15) {بررسی قاب\\پرچم، \lr{FEC / CRC}};
    \node[architecturebox] (phyout) at (4.375,-2.15) {دمدولاسیون\\تن، زمان‌بندی};
    \node[architecturebox] (air) at (1.325,-2.15) {کانال هوا\\بلندگو، هوا، میکروفون};
    \draw[arrow] (app)--(source);
    \draw[arrow] (source)--(packet);
    \draw[arrow] (packet)--(link);
    \draw[arrow] (link)--(phy);
    \draw[arrow] (phy)--(air);
    \draw[arrow] (air)--(phyout);
    \draw[arrow] (phyout)--(linkout);
    \draw[arrow] (linkout)--(sourceout);
    \draw[arrow] (sourceout)--(out);
\end{tikzpicture}
\end{otherlanguage}
\caption{مسیر دادۀ پیاده‌سازی‌شدۀ \lr{SonicDrop}. پردازش اختیاری گذرکلید پیش از کدگذاری پیوند بر محتوای بسته اعمال می‌شود؛ در دریافت، صدا پیش از رمزگشایی قاب فیلتر و جداسازی می‌شود.}
\label{fig:architecture}
\end{figure}

\subsection{اجزای مخزن کد}
\begin{table}[ht]
\centering
\small
\begin{tabularx}{\textwidth}{@{}Xl@{}}
\toprule
\textbf{وظیفه} & \textbf{مؤلفه} \\
\midrule
تبدیل \lr{UTF-8}/بیت، سریال‌سازی تصویر، \lr{RLE} و تبدیل بیتی اختیاری مبتنی بر گذرکلید. & \code{source\_coding.py} \\
سرآیند بسته، \lr{CRC-8}، همینگ \lr{(7,4)}، پرچم‌ها و درج بیت \lr{HDLC}، جریان‌های با پرچم مشترک و تزریق خطای کنترل‌شده. & \code{channel\_coding.py} \\
تولید شکل‌موج \lr{M-FSK}، سریال‌سازی \lr{WAV}، دمدولاسیون مبتنی بر انرژی تن، جست‌وجوی جریان و محاسبۀ نرخ‌های کانال. & \code{modem\_phy.py} \\
داشبورد \lr{Django}، نقطۀ پایانی ارسال/تحلیل، نقطۀ پایانی رمزگشایی \lr{WAV} و ثبت ارسال‌ها. & \code{views.py} \\
فرستندۀ واکنش‌گرا، گیرندۀ میکروفونی، بوم فکس، آزمایشگاه سیگنال، نمودار شکل‌موج و هماهنگی سمت مرورگر. & \code{index.html} \\
تاریخچۀ ارسال مبتنی بر \lr{SQLite} و معیارهای گزارش‌شدۀ پیوند. & \code{models.py} \\
\bottomrule
\end{tabularx}
\caption{ماژول‌های اصلی پیاده‌سازی.}
\label{tab:modules}
\end{table}

\section{عملکرد سرتاسری}
\label{sec:operation}

\Cref{fig:sequence} یک نشست معمول ارسال و دریافت را خلاصه می‌کند. فرستنده از نقطۀ پایانی \lr{Django} یک شکل‌موج کدگذاری‌شده می‌خواهد. نتیجه شامل بایت‌های \lr{WAV} و فراداده‌های تشخیصی است و مرورگر آن را پخش می‌کند. به‌طور مستقل، گیرنده می‌تواند صدا را ضبط کند یا یک فایل \lr{WAV} را رمزگشایی کند. رمزگشا تنظیمات مودم کاندیدا را آزمایش می‌کند، بسته‌ها را استخراج می‌کند و آخرین شناسۀ پیام را بازسازی می‌کند.

\begin{figure}[ht]
\centering
\begin{otherlanguage}{english}
\begin{tikzpicture}[>=Stealth]
    \node[draw=primaryblue,fill=lightblue,rounded corners=2pt,minimum width=3.2cm,
        minimum height=0.72cm,align=center,font=\small\bfseries] at (13.2,0) {\fa{مرورگر فرستنده}};
    \node[draw=primaryblue,fill=lightblue,rounded corners=2pt,minimum width=3.2cm,
        minimum height=0.72cm,align=center,font=\small\bfseries] at (8.8,0) {\fa{\lr{Django} / مودم}};
    \node[draw=primaryblue,fill=lightblue,rounded corners=2pt,minimum width=3.2cm,
        minimum height=0.72cm,align=center,font=\small\bfseries] at (4.4,0) {\fa{کانال صوتی}};
    \node[draw=primaryblue,fill=lightblue,rounded corners=2pt,minimum width=3.2cm,
        minimum height=0.72cm,align=center,font=\small\bfseries] at (0,0) {\fa{مرورگر گیرنده}};
    \draw[densely dashed,gray] (13.2,-0.4)--(13.2,-8.3);
    \draw[densely dashed,gray] (8.8,-0.4)--(8.8,-8.3);
    \draw[densely dashed,gray] (4.4,-0.4)--(4.4,-8.3);
    \draw[densely dashed,gray] (0,-0.4)--(0,-8.3);
    \draw[arrow] (13.2,-1.1)--node[fill=white,inner sep=2pt,font=\footnotesize]{\fa{\lr{POST}: بار داده + تنظیمات}}(8.8,-1.1);
    \draw[arrow] (8.8,-2.05)--node[fill=white,inner sep=2pt,font=\footnotesize]{\fa{\lr{WAV} + معیارها}}(13.2,-2.05);
    \draw[arrow] (13.2,-3.0)--node[fill=white,inner sep=2pt,font=\footnotesize]{\fa{بلندگو \lr{WAV} را پخش می‌کند}}(4.4,-3.0);
    \draw[arrow] (4.4,-3.95)--node[fill=white,inner sep=2pt,font=\footnotesize]{\fa{صدای ضبط‌شده}}(0,-3.95);
    \draw[arrow] (0,-4.9)--node[fill=white,inner sep=2pt,font=\footnotesize]{\fa{\lr{POST}: \lr{WAV} + تنظیمات}}(8.8,-4.9);
    \draw[arrow] (8.8,-5.85)--node[fill=white,inner sep=2pt,font=\footnotesize]{\fa{قاب‌های رمزگشایی‌شده + محتوا}}(0,-5.85);
    \draw[dashedarrow] (0,-6.8)--node[fill=white,inner sep=2pt,font=\footnotesize]{\fa{تن \lr{ACK/NACK} اختیاری}}(4.4,-6.8);
    \draw[dashedarrow] (4.4,-7.75)--node[fill=white,inner sep=2pt,font=\footnotesize]{\fa{فرستنده مدت کوتاهی گوش می‌دهد}}(13.2,-7.75);
    \node[fill=white,font=\scriptsize,align=center,text width=12.8cm] at (6.6,-9.0)
        {پیکان‌های خط‌چین نمایش بازخورد در مرورگر را نشان می‌دهند، نه یک پروتکل انتقال کاملاً مذاکره‌شده.};
\end{tikzpicture}
\end{otherlanguage}
\caption{دنبالۀ درخواست، صدا، رمزگشایی و بازخورد اختیاری.}
\label{fig:sequence}
\end{figure}

\subsection{نقطۀ پایانی ارسال}
نقطۀ پایانی \code{/api/transmit/} یک \lr{JSON} شامل نوع بار داده (\lr{payload})، بیت‌های متن یا فکس، باند حامل، مرتبۀ مدولاسیون، نرخ نماد، حالت تحویل و گزینه‌هایی برای \lr{RLE}، کدگذاری همینگ، نویز، تزریق خطا و استفاده از گذرکلید می‌پذیرد. این نقطۀ پایانی یک یا چند بستۀ بار داده می‌سازد، آن‌ها را قاب‌بندی می‌کند، کل جریان صوتی را تولید می‌کند و یک \lr{WAV} با کدگذاری \lr{base64} به‌همراه زمان‌بندی بسته‌ها و داده‌های تشخیصی برمی‌گرداند. تحلیل ارسال برای آزمایشگاه سیگنال نیز در فرایند سرور انجام می‌شود.

\subsection{نقطۀ پایانی دریافت}
نقطۀ پایانی \code{/api/decode-wav/} یک فایل \lr{WAV} و تنظیمات گیرنده را می‌پذیرد. این نقطۀ پایانی مقادیر نمونۀ \lr{PCM} تک‌کاناله را از جریان \lr{WAV} می‌خواند، فهرستی محدود از پیکربندی‌های مدولاسیون را امتحان می‌کند و نخستین پیکربندی‌ای را که قاب‌های معتبر تولید کند برمی‌گرداند. قاب‌ها با \lr{CRC} بررسی می‌شوند، سرآیندهای آن‌ها تجزیه می‌شود، پردازش اختیاری گذرکلید و رمزگشایی \lr{RLE} اعمال می‌شود و داده‌های متن یا فکس بر اساس شمارۀ ترتیب مونتاژ می‌شوند. اگر هیچ قاب معتبری یافت نشود، نقطۀ پایانی به‌جای داده‌های بازیابی‌شدۀ برنامه، خطا برمی‌گرداند.

\subsection{میکروفون زنده و هماهنگی مرورگر}
رابط کاربری از \lr{API}‌های صوتی مرورگر برای ضبط ورودی میکروفون استفاده می‌کند و ضبط‌های \lr{WAV} را به نقطۀ پایانی رمزگشایی می‌فرستد. یک کانال سمت مرورگر تنظیمات و رویدادهای نمایشی را بین زبانه‌ها همگام می‌کند. فرستنده همچنین می‌تواند پس از پخش، یک پنجرۀ کوتاه میکروفون باز کند تا به‌دنبال تن \lr{ACK/NACK} بگردد. از آنجا که زمان‌بندی زندۀ بسته‌ها در رابط کاربری و برخی کنترل‌های نمایشی با پیام‌های مرورگر هماهنگ می‌شوند، یک انیمیشن موفق یا بستۀ نمایش‌داده‌شده را نباید با رمزگشایی صوتیِ مستقل و موفق اشتباه گرفت؛ مسیری که بازیابی شکل‌موج ضبط‌شده را به کار می‌گیرد، نقطۀ پایانی رمزگشایی \lr{WAV} است.

\section{کدگذاری منبع و قالب بسته}
\label{sec:source}

\subsection{نمایش متن و تصویر}
متن به‌صورت \lr{UTF-8} کدگذاری و سپس با هشت بیت برای هر بایت نمایش داده می‌شود. رمزگشا بایت‌های ناقصِ انتهایی را حذف می‌کند و برای \lr{UTF-8} نامعتبر از نویسۀ جایگزین استفاده می‌کند؛ یک بررسی متن قابل‌چاپ، نتایج احتمالاً نویزی را فیلتر می‌کند. فکس $16\times16$ سطر‌به‌سطر به ۲۵۶ مقدار پیکسل دودویی مسطح می‌شود. هر پیکسل صفر/یک به‌صورت یک بیت سریال‌سازی می‌شود و گیرنده دنبالۀ بازیابی‌شده را دوباره به شبکه تبدیل می‌کند.

\subsection{کدگذاری طول‌دنباله (\lr{RLE})}
قالب \lr{RLE} هر دنباله (\lr{run}) را به‌صورت یک بیت برای مقدار تکرارشونده و به‌دنبال آن یک طول بدون‌علامت هفت‌بیتی نمایش می‌دهد. بنابراین هر رکورد هشت بیت است؛ طول دنباله‌ها حداکثر ۱۲۷ است و دنباله‌های بلندتر تقسیم می‌شوند. برای مثال، دنبالۀ پنج‌تایی از یک‌ها با رکورد \code{1 0000101} نمایش داده می‌شود. \lr{RLE} برای متن اختیاری است، در حالی که داده‌های فکس در پیاده‌سازی ارسال با \lr{RLE} کدگذاری می‌شوند. این روش برای دنباله‌های بلند (که در تصاویر دودویی ساده رایج است) مؤثر است، اما می‌تواند داده‌ای را که مکرراً تغییر می‌کند منبسط کند، زیرا هر دنباله هشت بیت کدگذاری‌شده مصرف می‌کند.

\subsection{سرآیند بسته}
هر بار داده با یک سرآیند ۱۲ بیتی آغاز می‌شود و پس از آن محتوای برنامه (یا محتوای تبدیل‌شده) می‌آید. \Cref{tab:header} فیلدها را دقیقاً شرح می‌دهد.

\begin{table}[ht]
\centering
\begin{tabularx}{\textwidth}{@{}Xlcc@{}}
\toprule
\textbf{معنی} & \textbf{فیلد} & \textbf{عرض} & \textbf{محدودۀ بیت} \\
\midrule
یک پیام را شناسایی می‌کند؛ مقادیر غیرصفر ۱ تا ۷ تولید می‌شوند. & شناسۀ پیام & 3 & ۰ تا ۲ \\
پرچم حالت جریانی/کلمه‌محور. & حالت تحویل & 1 & 3 \\
آخرین بستۀ پیام را مشخص می‌کند. & بستۀ آخر & 1 & 4 \\
پرچم فکس؛ صفر نشان‌دهندۀ متن است. & نوع بار داده & 1 & 5 \\
نشان می‌دهد که تبدیل گذرکلید اعمال شده است. & پرچم رمزگذاری & 1 & 6 \\
شمارۀ بسته که به‌پیمانۀ ۳۲ نمایش داده می‌شود. & شمارۀ ترتیب & 5 & ۷ تا ۱۱ \\
\bottomrule
\end{tabularx}
\caption{سرآیند بسته که ماژول کدگذاری کانال آن را بسته‌بندی می‌کند.}
\label{tab:header}
\end{table}

در حالت کلمه‌محور، متن به کلمات تقسیم و به‌صورت دنباله‌ای از بسته‌ها ارسال می‌شود؛ در حالت پیام کامل، حداکثر چهار کلمه در هر بسته گروه‌بندی می‌شود. محتوای فکس به قطعه‌های هم‌اندازه از بیت‌های فشرده تقسیم می‌شود (۳۲ بیت در حالت جریانی و ۶۴ بیت در حالت کامل). سرآیند بسته به گیرنده امکان می‌دهد داده‌ها را مرتب کند و بستۀ آخر را تشخیص دهد. فیلد پنج‌بیتیِ ترتیب، شماره‌گذاری بدون ابهام را به ۳۲ موقعیت بسته محدود می‌کند؛ این موضوع برای بدترین الگوهای فکس که ضعیف فشرده می‌شوند اهمیت دارد.

\subsection{تبدیل اختیاری گذرکلید و محدودیت‌های امنیتی}
وقتی گذرکلید غیرخالی ارائه شود، پیاده‌سازی با درهم‌سازی (\lr{hash}) رشته‌هایی شامل برچسب نسخه، گذرکلید، شمارۀ ترتیب بسته (به‌عنوان نمک) و یک شمارنده با \lr{SHA-256}، یک جریان بیتی تکرارپذیر استخراج می‌کند. این جریان با بار داده‌ای که یک برچسب راستی‌آزمایی هشت‌بیتیِ مشتق‌شده از \lr{SHA-256} پیش از آن قرار گرفته، \lr{XOR} می‌شود. گیرنده جریان را بازتولید می‌کند، نتیجه را \lr{XOR} می‌کند و برچسب را مقایسه می‌کند. شمارۀ ترتیب بسته به‌عنوان نمک استفاده می‌شود تا بسته‌های مختلفِ یک پیام معمولاً ورودی‌های تبدیل متفاوتی داشته باشند.

\textbf{هشدار امنیتی:} این یک نمایش پروژه‌ای است، نه یک طرح رمزنگاری یا احراز هویت بازبینی‌شده. این ساختار سفارشی فاقد تابع استخراج کلید مقاوم در برابر حدس گذرواژه (\lr{KDF}) است، از یک برچسب راستی‌آزمایی بسیار کوتاه (۸ بیتی) استفاده می‌کند، هیچ تضمین استاندارد \lr{AEAD} ندارد و در صورت تکرار همان کلید و همان شمارۀ ترتیب بسته، ورودی‌های جریان کلید می‌توانند دوباره استفاده شوند. از آن نباید برای محافظت از اطلاعات حساس استفاده شود. یک سامانۀ تولیدی باید از یک ساختار رمزگذاری احراز‌هویت‌شدۀ بازبینی‌شده همراه با مدیریت \lr{nonce} و، در صورت لزوم، یک تابع استخراج کلید مبتنی بر گذرواژه استفاده کند.

\section{قاب‌بندی و کنترل خطا در لایۀ پیوند داده}
\label{sec:link}

\subsection{\lr{CRC-8}}
لایۀ پیوند یک باقی‌ماندۀ \lr{CRC} هشت‌بیتی را که با چندجمله‌ای مولد زیر محاسبه شده است ضمیمه می‌کند (الگوی بیتی \lr{100000111}):
\[
    G(x)=x^8+x^2+x+1 .
\]
\lr{CRC} روی سرآیند و محتوای بسته، پیش از \lr{FEC} محاسبه می‌شود. در گیرنده، باقی‌ماندۀ صفر نشان می‌دهد داده‌های محافظت‌شدۀ دریافتی بررسی \lr{CRC} را می‌گذرانند. \lr{CRC} بسیاری از خرابی‌های تصادفی را تشخیص می‌دهد؛ این کد، کد تصحیح خطا نیست و نباید آن را احراز هویت رمزنگاری تلقی کرد.

\subsection{تصحیح خطای پیشرو با کد همینگ \lr{(7,4)}}
در صورت فعال بودن، سرآیند ۱۲ بیتی به‌همراه محتوا و \lr{CRC} تا مضربی از چهار بیت پدینگ می‌شود و در بلوک‌های چهاربیتی داده کدگذاری می‌شود. سه بیت توازن در موقعیت‌های ۱، ۲ و ۴ قرار می‌گیرند و برای هر چهار بیت ورودی، هفت بیت ارسال می‌شود. رمزگشا یک سندرم سه‌بیتی محاسبه می‌کند؛ سندرم ناصفر موقعیت بیتی را که باید معکوس شود مشخص می‌کند. این کد یک بیت را در هر کدواژۀ هفت‌بیتی تصحیح می‌کند. این روش خطاهای چندگانه در یک کدواژه را به‌طور قابل‌اعتماد تصحیح نمی‌کند و چنین خطاهایی می‌توانند به تصحیح نادرست بینجامند؛ \lr{CRC} پس از رمزگشایی بررسی یکپارچگیِ دوم را فراهم می‌کند.

نرخ کد پیش از سربار \lr{CRC}، قاب‌بندی و درج بیت برابر $4/7\approx0{,}571$ است. افزونگی اضافی زمان ارسال را افزایش می‌دهد، اما می‌تواند خطاهای بیتی منفردِ کانال را بدون درخواست نسخۀ دیگر بازیابی کند.

\subsection{مرزهای قاب به سبک \lr{HDLC}}
کدکنندۀ کانال یک قاب را از یک مقدمۀ متناوب ۲۴ بیتی، پرچم \code{01111110}، یک بدنۀ محافظت‌شده و یک پرچم پایانی می‌سازد. در بدنه، درج بیت پس از هر دنبالۀ پنج‌تایی از یک‌ها یک صفر درج می‌کند. رمزگشا بیت‌های درج‌شده را پیش از پردازش \lr{FEC} و \lr{CRC} حذف می‌کند. این کار احتمال اینکه داده‌های بدنه توالی پرچم را تقلید کنند کاهش می‌دهد. برای چند بسته، جریان مشترک شامل پرچم‌های تکراری میان بدنه‌ها برای جداسازی بیشتر، و سپس یک قطعۀ مقدمۀ پایانی کوتاه است.

\begin{figure}[ht]
\centering
\begin{otherlanguage}{english}
\begin{tikzpicture}[>=Stealth]
    \tikzset{framefield/.style={draw=primaryblue,rounded corners=2pt,
        minimum width=3.15cm,minimum height=1.6cm,text width=2.95cm,
        align=center,font=\footnotesize}}
    \node[framefield,fill=lightblue] at (11.625,0)
        {\textbf{مقدمه}\\[2pt]۲۴ بیت متناوب};
    \node[framefield,fill=lightgray] at (8.275,0)
        {\textbf{جداکنندۀ آغاز}\\[2pt]پرچم \lr{HDLC}: ۸ بیت};
    \node[framefield,fill=teal!14] at (4.925,0)
        {\textbf{بدنۀ محافظت‌شده}\\[2pt]سرآیند / داده / \lr{CRC-8}};
    \node[framefield,fill=lightgray] at (1.575,0)
        {\textbf{جداکنندۀ پایان}\\[2pt]پرچم \lr{HDLC}: ۸ بیت};
    \node[font=\small\bfseries] at (6.6,-1.75) {پردازش اعمال‌شده بر بدنۀ قاب};
    \tikzset{protectionstep/.style={draw=darkgray,rounded corners=2pt,
        minimum width=4.1cm,minimum height=0.85cm,text width=3.8cm,
        align=center,font=\small}}
    \node[protectionstep,fill=lightblue] (crc) at (11.15,-2.6) {افزودن \lr{CRC-8}};
    \node[protectionstep,fill=teal!14] (ham) at (6.7,-2.6) {همینگ \lr{(7,4)} اختیاری};
    \node[protectionstep,fill=orange!15] (stuff) at (2.25,-2.6) {درج بیت \lr{HDLC}};
    \draw[arrow] (crc.west)--(ham.east);
    \draw[arrow] (ham.west)--(stuff.east);
    \node[font=\footnotesize,align=center,text width=13.2cm] at (6.6,-3.8)
        {در جریان چندبسته‌ای، پرچم‌های تکراری میان بدنه‌ها درج می‌شوند؛ حفاظت بدنه پس از مونتاژ سرآیند و محتوای بسته اعمال می‌شود.};
\end{tikzpicture}
\end{otherlanguage}
\caption{فیلدهای قاب با عرض برابر و ترتیب حفاظت بدنه. عرض فیلدها نمایشی است؛ تنها مقدمه و پرچم‌ها طول بیتی ثابت دارند.}
\label{fig:frame}
\end{figure}

\subsection{تحویل بسته و نمایش ارسال مجدد}
رابط کاربری قابلیت بازخورد خود را «\lr{ARQ} ترکیبی» (\lr{Hybrid ARQ}) نام‌گذاری کرده است. مرورگر می‌تواند یک تن \lr{ACK} با فرکانس \lr{6.8} کیلوهرتز یا یک تن \lr{NACK} با فرکانس \lr{3.4} کیلوهرتز تولید کند و فرستنده می‌تواند پس از پخش به این فرکانس‌ها گوش دهد. گزینۀ \lr{NACK} اجباری یک مسیر نمایشی فراهم می‌کند که یک بستۀ انتخابی را خراب می‌کند و بسته را بر اساس شمارۀ ترتیب دوباره ارسال می‌کند؛ فرستنده در آن جریان رابط کاربری حداکثر یک تلاش مجدد خودکار را مجاز می‌داند. نقطۀ پایانی رمزگشا نیز وقتی قاب‌های معتبر بازیابی کرده باشد، فراداده‌های مرتبط با \lr{ACK} را برمی‌گرداند.

این هنوز یک پروتکل کاملِ \lr{ARQ} با سازوکار انتخاب و تکرار (\lr{selective-repeat}) نیست. این پروتکل یک بستۀ بازخورد پایدار روی هوا با شناسه‌های پیام و شمارۀ ترتیب، زمان‌سنج‌های مذاکره‌شده، سقف تلاش مجدد در سطح نشست، تضمین حذف تکراری‌ها یا بازیابی بسته‌های گمشده برای همۀ مسیرهای دریافت ارائه نمی‌دهد. به‌ویژه، شبیه‌سازی محلی مرورگر، تحلیل \lr{API} و دریافت مستقل \lr{WAV} سطوح متفاوتی از اتصال دارند. بنابراین این گزارش \lr{ARQ} را یک نمایش بازخوردِ نمونۀ اولیه توصیف می‌کند، نه تضمینی برای عدم از دست رفتن داده.

\section{لایۀ فیزیکی}
\label{sec:phy}

\subsection{انتخاب‌های مدولاسیون}
لایۀ فیزیکی از مدولاسیون کلیدزنی جابه‌جایی فرکانس $M$-تایی (\lr{M-ary FSK}) استفاده می‌کند. هر نماد یک حامل را انتخاب می‌کند؛ برای $M=2$ هر نماد یک بیت و برای $M=4$ هر نماد دو بیت حمل می‌کند. نگاشت نماد در \lr{4-FSK} مقادیر \code{00}، \code{01}، \code{10} و \code{11} را به‌ترتیب به چهار فرکانس حاملِ پیکربندی‌شده اختصاص می‌دهد. اگر ورودی \lr{4-FSK} تعداد فرد بیت داشته باشد، مدولاتور یک بیت صفر پدینگ می‌کند.

\begin{table}[ht]
\centering
\begin{tabularx}{\textwidth}{@{}ccXcl@{}}
\toprule
\textbf{بیت/نماد} & \textbf{پهنای باند مصرفی} & \textbf{فرکانس‌های حامل (هرتز)} & \textbf{\lr{M}} & \textbf{حالت}\\
\midrule
1 & 2000 هرتز & 2600، 3800 & 2 & شنیداری\\
2 & 2600 هرتز & 2400، 3100، 3800، 4500 & 4 & شنیداری\\
1 & 2000 هرتز & 14800، 15800 & 2 & نزدیک‌به‌فراصوت\\
2 & 2400 هرتز & 14400، 15000، 15600، 16200 & 4 & نزدیک‌به‌فراصوت\\
\bottomrule
\end{tabularx}
\caption{مجموعه‌های حامل و مقادیر پهنای باند پیکربندی‌شده در لایۀ فیزیکی.}
\label{tab:tones}
\end{table}

در نرخ نمونه‌برداری $f_s$ و نرخ نماد $R_s$، پیاده‌سازی از
\[
    N_s=\left\lfloor\frac{f_s}{R_s}\right\rfloor
\]
نمونه در هر نماد استفاده می‌کند. نرخ نمونه‌برداری پیش‌فرض \lr{44.1} کیلوهرتز و نرخ نماد پیش‌فرض ۴۵ بود است. نرخ بیت اسمیِ بدون کدگذاریِ حاصل برابر است با
\[
    R_b=R_s\log_2(M).
\]
بنابراین در ۴۵ بود، \lr{2-FSK} ۴۵ بیت بر ثانیۀ اسمی و \lr{4-FSK} ۹۰ بیت بر ثانیۀ اسمی حمل می‌کند، پیش از قاب‌بندی، \lr{CRC}، \lr{FEC}، درج بیت و فشرده‌سازی اختیاری. این‌ها نرخ‌های برگرفته از پیکربندی هستند، نه گذردهیِ مفیدِ (\lr{goodput}) اندازه‌گیری‌شدۀ برنامه.

\subsection{تولید شکل‌موج}
مدولاتور برای هر نماد یک سینوسی در حاملِ انتخاب‌شده تولید می‌کند. فاز را از نمادی به نماد بعدی حفظ می‌کند و یک شیب کوتاه کسینوس‌بالابرده (\lr{raised-cosine}) در هر لبۀ نماد اعمال می‌کند تا گذارهای ناگهانی دامنه کاهش یابد. شکل‌موج اعشاریِ حاصل مقیاس‌بندی و به \lr{WAV} تک‌کانالۀ \lr{PCM} شانزده‌بیتیِ علامت‌دار تبدیل می‌شود. برای \lr{4-FSK}، یک بیت انتهاییِ فرد، همان‌طور که اشاره شد، پدینگ می‌شود.

\begin{figure}[ht]
\centering
\begin{otherlanguage}{english}
\begin{tikzpicture}
\begin{axis}[
    width=0.9\textwidth,
    height=5.0cm,
    xlabel={\lr{Symbol rate $R_s$ (baud)}},
    ylabel={\lr{Uncoded bit rate $R_b$ (bit/s)}},
    xmin=20,xmax=95,ymin=0,ymax=205,
    xtick={25,30,45,65,90},
    ytick={0,30,45,60,90,130,180},
    grid=both,
    legend style={at={(0.02,0.98)},anchor=north west,draw=none,fill=white},
    axis line style={darkgray},
    tick label style={font=\small},
    label style={font=\small}
]
    \addplot[primaryblue,thick,domain=25:90,samples=2] {x};
    \addlegendentry{\lr{2-FSK: $R_b=R_s$}}
    \addplot[teal,thick,domain=25:90,samples=2] {2*x};
    \addlegendentry{\lr{4-FSK: $R_b=2R_s$}}
    \addplot[only marks,mark=*,mark size=2pt,orange] coordinates {(45,45) (45,90)};
\end{axis}
\end{tikzpicture}
\end{otherlanguage}
\caption{نرخ بیت اسمی مدولاسیون در برابر نرخ نماد، محاسبه‌شده از $R_b=R_s\log_2(M)$. نشانگرها تنظیم پیش‌فرض ۴۵ بود پروژه را نشان می‌دهند. این نمودار سربار پروتکل را شامل نمی‌شود و نمایانگر گذردهی اندازه‌گیری‌شده نیست.}
\label{fig:rate}
\end{figure}

\subsection{دمدولاسیون و جست‌وجوی زمان‌بندی}
برای هر نماد، گیرنده ناحیۀ محافظ لبه (تقریباً \lr{14\%} نماد در هر انتها) را کنار می‌گذارد و نمونه‌های باقی‌مانده را با الگوهای هم‌فاز و متعامد (\lr{in-phase / quadrature}) هر حامل کاندیدا همبسته می‌کند. مجذور اندازۀ همبستگی‌ها انرژی تن را تخمین می‌زند. وقتی تعداد کافی نماد وجود دارد، انرژی هر حامل بر یک صدک بالای توان مشاهده‌شدۀ آن نرمال می‌شود تا پاسخ فرکانسی نامساویِ بلندگو/میکروفون تا حدی جبران شود. حاملی که بیشترین انرژی نرمال‌شده را دارد انتخاب می‌شود؛ نمادهای کم‌انرژی حذف (\lr{squelch}) می‌شوند.

برای دریافت \lr{WAV}، پیاده‌سازی میانگین ضبط را حذف می‌کند، یک ماسک میان‌گذر \lr{FFT} پیرامون بازۀ حامل پیکربندی‌شده اعمال می‌کند، بهره را نرمال می‌کند و با استفاده از غالب بودن تن، یک بازۀ فعال (\lr{burst}) را جدا می‌کند. سپس تا ۲۰ انحراف را در یک بازۀ نماد امتحان می‌کند. در هر انحراف، بیت‌های حاصل برای یافتن پرچم قاب پویش و با \lr{FEC} و \lr{CRC} بررسی می‌شوند. این جست‌وجو به تحمل نقطۀ شروع نامعلوم کمک می‌کند، اما یک جست‌وجوی زمان‌بندیِ محدود است، نه یک حلقۀ بازیابی کامل زمان‌بندی یا یک هم‌ساز کانال (\lr{equalizer}).

\subsection{کار در باند شنیداری در برابر نزدیک‌به‌فراصوت}
مجموعۀ حامل نزدیک‌به‌فراصوت می‌کوشد تن‌های شنیدنی را کاهش دهد و در عین حال در حد نایکوئیست اسمیِ نرخ نمونه‌برداری \lr{44.1} کیلوهرتز بماند. سخت‌افزار واقعی ممکن است این فرکانس‌ها را به‌شدت تضعیف کند یا از طریق پردازش صوتی داخلی آن‌ها را حذف کند؛ برخی میکروفون‌ها، بلندگوها و کدک‌ها آن‌ها را به‌اندازۀ کافی حفظ نمی‌کنند. تنظیم شنیداری معمولاً مشاهده و اشکال‌زدایی را آسان‌تر می‌کند. نباید فرض کرد که هیچ‌یک از حالت‌های حامل بدون اندازه‌گیری روی همۀ دستگاه‌ها کار می‌کند.

\section{معیارهای کانال و تفسیر آن‌ها}
\label{sec:metrics}

لایۀ فیزیکی یک \lr{SNR} تخمینی بر حسب دسی‌بل، یک نرخ بیت اسمی، یک نرخ نایکوئیست و یک تخمین ظرفیت شانون گزارش می‌کند. در کد، این تخمین‌ها از توان تن آشکارشده و توان‌های باقی‌ماندۀ تن‌های رقیب استفاده می‌کنند:
\[
    \mathrm{SNR}_{dB}=10\log_{10}\!\left(\frac{P_s}{P_n}\right),
\]
که در آن $P_s$ میانگین توان تخمینیِ حاملِ برنده و $P_n$ میانگین توان در سایر فرضیه‌های حامل است. حدهای نظریِ گزارش‌شده عبارت‌اند از
\[
    R_{\mathrm{Nyquist}}=2B\log_2(M), \qquad
    C_{\mathrm{Shannon}}=B\log_2(1+\mathrm{SNR}_{linear}),
\]
با پهنای باند پیکربندی‌شدۀ $B$. این‌ها محاسبات مرجع تحت مفروضات آرمانی هستند؛ تضمینی نیست که کانال صوتی این مفروضات را برآورده کند و \lr{SNR} تخمینی، اندازه‌گیری آزمایشگاهی احتمال خطای بیت سرتاسری نیست.

نقطۀ پایانی ارسال و تحلیل یک \lr{WAV} تمیز و سنتز‌شده برمی‌گرداند، سپس نویز گاوسی و در صورت تمایل وارونگی‌های بیتی مصنوعی را روی یک نسخۀ جداگانۀ قاب نخست برای تحلیل اعمال می‌کند. \lr{BER} آن با مقایسۀ آن نسخۀ دمدوله‌شده با قاب تولیدشده محاسبه می‌شود. این معیار برای آزمایش‌های کنترل‌شده مفید است، اما اندازه‌گیری \lr{BER} میدانی نیست. در مسیر دریافت \lr{WAV}، مقدار شبیه‌\lr{BER} گزارش‌شده از بیت‌های تصحیح‌شده و طول قابِ نمایش‌داده‌شده به دست می‌آید؛ نباید آن را برآورد آماری کاملی روی بسته‌های مستقلِ بسیار تفسیر کرد.

\section{برنامۀ وب و پایداری داده}
\label{sec:web}

\lr{Django} سه مسیر ارائه می‌دهد: داشبورد در \code{/}، تحلیل ارسال در \code{/api/transmit/} و رمزگشایی \lr{WAV} در \code{/api/decode-wav/}. \lr{NumPy} عملیات شکل‌موج و طیفی را پیاده‌سازی می‌کند. داشبورد ورودی متن و فکس، کنترل‌هایی برای حامل و مدولاسیون و گزینه‌های پیوند، یک گیرندۀ ضبط، بارگذاری \lr{WAV}، بازرس بیت قاب، کارت‌های معیار و یک نمودار شکل‌موج فراهم می‌کند. شبکۀ فکس $16\times16$ پیکسل است.

هر درخواست ارسال و تحلیل یک رکورد \code{TransmissionLog} در \lr{SQLite} ایجاد می‌کند که شامل باند انتخابی، فیلد مرتبط با نرخ نماد، گزینه‌های کدگذاری، طول بیت‌ها، تعداد خطاهای تزریق‌شده و تصحیح‌شده، \lr{BER}، \lr{SNR}، نرخ‌های نظری، نتیجۀ \lr{CRC} و برچسب بازیابی‌شده یا نمایشی است. داشبورد ورودی‌های اخیر را نشان می‌دهد. این تاریخچه از آزمایش‌ها پشتیبانی می‌کند، اما یک سامانۀ تله‌متری یا حساب کاربری تولیدی نیست.

\section{آزمون و ارزیابی}
\label{sec:testing}

مخزن کد شامل یک نمایش موتور است که:
\begin{enumerate}
    \item یک جملۀ آزمون \lr{UTF-8} را به بیت تبدیل می‌کند، آن را با \lr{CRC-8} و کدگذاری همینگ قاب‌بندی می‌کند، خطاهای مصنوعی تزریق می‌کند و نتایج تصحیح رمزگشا و \lr{CRC} را چاپ می‌کند.
    \item یک تصویر دودویی نمونه را با \lr{RLE} کدگذاری و رمزگشایی می‌کند و تأیید می‌کند که بیت‌های خروجی با بیت‌های ورودی مطابقت دارند.
\end{enumerate}
فایل \lr{README} دستور \code{manage.py check} در \lr{Django} و یک نمایش اختیاری بلندگو/میکروفون را مستند می‌کند. این بررسی‌ها یک محک (\lr{benchmark}) بین‌دستگاهی نیستند؛ هیچ تضمینی دربارۀ برد، نرخ خطا، امنیت یا سازگاری سخت‌افزاری ادعا نمی‌شود.

\subsection{بررسی‌های محلیِ تکرارپذیر}
با نصب وابستگی‌های پروژه، بررسی‌های مستندشده عبارت‌اند از:
\begin{otherlanguage}{english}
\begin{lstlisting}[basicstyle=\ttfamily\small,frame=single,backgroundcolor=\color{lightgray}]
python manage.py check
python transceiver/engine/test_engine.py
\end{lstlisting}
\end{otherlanguage}
برای ارزیابی صوتی سرتاسری، \lr{WAV} منتشرشده را با دستگاه دوم و با ثبت فاصله، جهت‌گیری، حجم صدا و شرایط محیطی ضبط کنید. هر شرایط را روی بسته‌های بسیار تکرار کنید، هم نرخ موفقیت بسته و هم گذردهی مفید را گزارش دهید و ضبط‌های اصلی را نگه دارید تا نتایج رمزگشایی قابل بازتولید باشند.

\section{محدودیت‌ها و کارهای آینده}
\label{sec:limitations}

\subsection{محدودیت‌های فعلی}
\begin{itemize}
    \item \textbf{رمزنگاری نمونۀ اولیه:} تبدیل \lr{SHA-256/XOR} و برچسب کوتاه برای داده‌های محرمانه ایمن نیست؛ برای هر نیاز امنیتی واقعی از یک کتابخانۀ استاندارد رمزگذاریِ احراز‌هویت‌شده استفاده کنید.
    \item \textbf{\lr{ARQ} ناقص:} بازخورد و رفتار تلاش مجدد عمدتاً یک نمایش مرورگری است. پروتکل کاملِ روی هوا برای شناسه‌ها، سیاست زمان‌سنجی، بازیابی از دست‌رفتگی و مدیریت تکراری‌ها را ندارد.
    \item \textbf{محدودیت اندازۀ بسته و ترتیب:} سرآیند فقط پنج بیت ترتیب دارد. پیکسل‌های فکس با تناوب زیاد می‌توانند رکوردهای \lr{RLE} کافی تولید کنند که از ۳۲ موقعیت بسته فراتر رود.
    \item \textbf{مدل زمان‌بندی و کانال:} گیرنده از یک جست‌وجوی فاز محدود، آشکارساز انرژی و جبران طیفی ساده استفاده می‌کند. پاسخ ضربۀ کامل را تخمین نمی‌زند و رانش ساعت را به‌طور پیوسته دنبال نمی‌کند.
    \item \textbf{تطبیق نرخ:} کد یک نمایۀ پیشنهادی از آستانه‌های \lr{SNR/BER} تولید می‌کند؛ این پیشنهاد تشخیصی است و به‌تنهایی تغییر خودکارِ مذاکره‌شده در یک پیوند در حال اجرا نیست.
    \item \textbf{اندازه‌گیری‌های محدود:} تزریق خطای مصنوعی و نمایش موتور موجود، برد صوتی، پایداری در اتاق‌های مختلف یا گذردهی واقعی را اثبات نمی‌کنند.
    \item \textbf{نام‌گذاری فیلد تاریخچه:} نمای ارسال در حال حاضر نرخ بیت مؤثر را در فیلد پایگاه دادۀ \code{baud\_rate} ذخیره می‌کند؛ مقادیر تاریخی را نباید اندازه‌گیری نرخ نماد دانست.
    \item \textbf{اعتبارسنجی ورودی و استقرار:} این یک برنامۀ توسعۀ محلی است. پیش از استقرار عمومی، محدودیت اندازۀ درخواست را اعمال کنید، همۀ گزینه‌ها را اعتبارسنجی کنید، سیاست \lr{CSRF/API} را بازبینی کنید، رازهای تولیدی و میزبان‌های مجاز را پیکربندی کنید و از پیکربندی مناسب سرور استفاده کنید.
\end{itemize}

\subsection{گام‌های بعدیِ پیشنهادی}
\begin{enumerate}
    \item آزمون‌های واحد خودکار برای بردارهای \lr{CRC}، قاب‌بندی/حذف درج بیت، حالت‌های همینگِ تک‌بیتی و چندبیتی، حالت‌های مرزی \lr{RLE}، مرزهای بسته، رفتار شکست رمزگذاری و رفت‌وبرگشت مودم اضافه کنید.
    \item یک قالب قاب نسخه‌دار با طول بار داده، شناسه‌های پایدار پیام و شمارۀ ترتیبِ به‌اندازۀ کافی بزرگ برای حداکثر اندازۀ پیام پشتیبانی‌شده اضافه کنید.
    \item تبدیل سفارشی گذرکلید را با یک طرح استاندارد \lr{AEAD} جایگزین کنید و استخراج کلید امن و چرخۀ عمر \lr{nonce} را تعریف کنید.
    \item یک قاب بازخورد واقعی یا دنبالۀ تن کدگذاری‌شدۀ مقاوم، شامل هویت بسته و قواعد زمان‌سنجی، تلاش مجدد و تکراری، مشخص کنید و روی دستگاه‌های مستقل آزمایش کنید.
    \item صدای ضبط‌شده را در دستگاه‌ها و محیط‌های مختلف اندازه‌گیری کنید. نرخ نمونه‌برداری، فاصله، حجم صدا، مجموعۀ حامل، نرخ بود، اندازۀ بسته، تعداد آزمایش‌ها، تعداد خطاها و فاصله‌های اطمینان را منتشر کنید.
    \item همگام‌سازی و تخمین کانال را بهبود دهید و درهم‌ریزی (\lr{interleaving}) یا کدهای قوی‌تر را برای نویز انفجاری و انعکاس‌های طولانی (\lr{reverberation}) ارزیابی کنید.
    \item معیارهای تولیدشده توسط شبیه‌ساز را از اندازه‌گیری‌های روی ضبط‌ها در رابط کاربری و تاریخچۀ ذخیره‌شده جدا کنید.
\end{enumerate}

\section{نتیجه‌گیری}
\label{sec:conclusion}

\lr{SonicDrop} یک مودم صوتی آموزشیِ سرتاسری است: داده‌های برنامه را به جریان‌های بیتیِ قاب‌بندی‌شده تبدیل می‌کند، این جریان‌ها را به تن‌های صوتی نگاشت می‌کند و یک گیرندۀ نرم‌افزاری برای بازیابی محتوا از ضبط‌ها فراهم می‌کند. پیاده‌سازی لایه‌ایِ آن، کدگذاری منبع، بسته‌سازی، \lr{CRC}، \lr{FEC} اختیاری، \lr{M-FSK}، معیارهای سیگنال و ارائۀ وب را قابل‌مشاهده و قابل‌آزمون می‌کند. شکل‌ها و نمودار نرخِ گنجانده‌شده الگوریتم‌های پیاده‌سازی‌شده و محاسبات نظری را توصیف می‌کنند، نه اینکه نتیجۀ محک ادعا کنند. این پروژه پایۀ مفیدی برای آزمایش‌های بیشتر است و رمزنگاری، کامل بودن پروتکل، ارزیابی صوتی و همگام‌سازی قوی‌تر به‌عنوان مهم‌ترین گام‌های بعدی شناسایی شده‌اند.

\appendix
\renewcommand{\thesection}{الف}
\section{مرجع سریع پیکربندی پیش‌فرض}
\label{app:defaults}
\begin{table}[ht]
\centering
\begin{tabularx}{\textwidth}{@{}Xl@{}}
\toprule
\textbf{پیش‌فرض / رفتار} & \textbf{پارامتر} \\
\midrule
\lr{44,100} نمونه بر ثانیه؛ خروجی \lr{WAV} تک‌کاناله با \lr{PCM} شانزده‌بیتی. & نرخ نمونه‌برداری \\
۴۵ بود در پیش‌فرض‌های \lr{API} و کنترل اولیۀ رابط؛ لغزندۀ رابط بازۀ ۲۵ تا ۹۰ بود را با گام‌های ۵ بودی می‌پوشاند. & نرخ نماد \\
به‌طور پیش‌فرض \lr{4-FSK} (دو بیت در هر نماد)؛ \lr{2-FSK} نیز قابل انتخاب است. & مدولاسیون \\
به‌طور پیش‌فرض شنیداری؛ نزدیک‌به‌فراصوت قابل انتخاب است. & باند حامل \\
به‌طور پیش‌فرض فعال؛ می‌توان آن را غیرفعال کرد. & همینگ \\
به‌طور پیش‌فرض غیرفعال؛ می‌توان آن را فعال کرد. محتوای فکس همیشه توسط خط‌لولۀ ارسال با \lr{RLE} کدگذاری می‌شود. & \lr{RLE} متن \\
نقطۀ پایانی تحلیل، نویز گاوسیِ پیکربندی‌پذیر و وارونگی‌های بیتیِ مصنوعیِ اختیاریِ تکی، چندگانه یا انفجاری را پشتیبانی می‌کند. & کنترل‌های نویز/خطا \\
تنظیمات درخواستی و چند پیکربندی جایگزین را امتحان می‌کند، سپس نخستین مجموعۀ قاب‌های معتبرِ یافت‌شده را برمی‌گرداند. & جست‌وجوی دریافت \lr{WAV} \\
\bottomrule
\end{tabularx}
\caption{مقادیر پیش‌فرض و رفتار قابل‌مشاهده در پیاده‌سازی فعلی.}
\end{table}

\end{document}